\sect{Датасеты партитур и обучающие данные}{datasets}
%
Качество отслеживания партитуры в реальном времени критически зависит
от того, насколько обучающая выборка покрывает разнообразие реального
исполнения: вариации темпа, ошибочные ноты, орнаменты и шум входного
канала. В настоящей работе использованы три источника данных,
дополняющих друг друга по охвату и контролируемости.

В качестве синтетического источника применены прелюдии Шопена~Op.~28,
переведённые из формата MusicXML в MIDI с помощью библиотеки
\textit{music21}. На основе каждой прелюдии строилось семейство из
сорока модифицированных исполнений: к нотным онсетам добавлялся
гауссовский шум $\sigma_\tau$ от $2$ до $12$\,\%, локальные
ритардандо и ачелерандо моделировались кусочно-линейным умножением
длительностей в окне восьми тактов, в случайных позициях вставлялись
пропуски и однократные повторы нот. Этот источник позволяет получать
ground-truth соответствие между партитурой и исполнением, что
необходимо для обучения с учителем.

Открытые корпуса MAPS~\cite{Mueller2007} и MSMD~\cite{Dorfer2018}
использованы для оценки эмиссионной модели в аудиорежиме и для
проверки переноса моделей между фортепианными произведениями
различных эпох. MAPS содержит около $270$ записей фортепианной музыки
с поточечной MIDI-разметкой, MSMD~--- $497$ произведений с попарным
выравниванием аудио и нотного изображения. Для соответствия задачам
конкретного концерта Чайковского из числа MAPS отобраны записи в
тональностях, близких к репертуару пилота.

Собственный корпус \textit{CU-Concerto-2026} собран в стенах
Центрального Университета: четыре пианиста-консерваторца исполнили
первые части фортепианных концертов Чайковского~b-moll и
Рахманинова~c-moll суммарной продолжительностью $42$ часа. Запись
велась в режиме MIDI через цифровое пианино Yamaha~CLP-735, что
обеспечило точные онсеты и значения velocity, а параллельная
аудиозапись на конденсаторный микрофон применялась для
тестирования аудиоветви системы. Каждое исполнение было размечено
вручную: расставлены границы такта, отмечены ферматы, ритардандо,
повторы, пропуски и фальшивые ноты. Сводная статистика приведена в
\Tabref{datasets-summary}.

Особое внимание уделено сохранению этики записи: участники подписали
информированное согласие на использование их исполнений в
исследовательских целях, исходные аудиофайлы хранятся локально на
сервере университета без публикации в открытом доступе. Производные
признаки (выровненные MIDI, тембральные признаки) предполагается
опубликовать как открытый датасет в составе сопровождающего материала
после защиты статьи.

\par\addvspace{2pt}
{\centering
\refstepcounter{table}\tablab{datasets-summary}%
\footnotesize
\begin{tabular}{@{}lcccc@{}}
  \toprule
  Датасет & Часов & Произв. & Темп.\,откл. & Ошиб. \\
  \midrule
  Шопен-синт.       & 28 & 24 & до 12\,\%   & контролируемые \\
  MAPS              & 18 & 64 & до 8\,\%    & нет \\
  MSMD              & 26 & 71 & до 6\,\%    & нет \\
  CU-Concerto-2026  & 42 & 12 & до 15\,\%   & естественные \\
  \bottomrule
\end{tabular}\par
\vspace{2pt}
\tablename~\thetable. Сводка использованных источников данных. Поле
«Темп.\,откл.»~--- максимальное относительное отклонение темпа от
номинального, «Ошиб.»~--- характер ошибок в датасете.\par}
\addvspace{4pt}

Разбиение выборки на обучающую, валидационную и тестовую части
проводилось по уровню \textit{произведений}, а не по уровню записей,
чтобы исключить утечку информации между тембрами одного и того же
концерта. Обучающая выборка использовалась для подбора параметров
эмиссионной модели и обучения свёрточного фронтенда (см.
\Secref{cnn-corner}); валидационная~--- для подбора гиперпараметров
$p_{\text{stay}}, p_{\text{advance}}, p_{\text{skip}}, R_{\max}$
гибридной схемы и для остановки обучения RL-модуля; тестовая~--- для
финальных метрик \Secref{experiments}.
